Este artigo investiga como os determinantes do retorno econômico esperado e o ambiente institucional condicionam a conversão de vegetação nativa para usos agropecuários em Mato Grosso. A análise parte da hipótese de que maiores retornos esperados elevam a pressão de conversão, com respostas distintas conforme a capacidade produtiva local e as condições de implementação das regras ambientais. Para examinar essas relações, propõe-se um painel territorial anual, tendo 2001--2024 como período candidato, que combine transições de uso da terra, preços agropecuários, capacidade produtiva predeterminada e informações sobre fiscalização e sanções. A estratégia econométrica explora a exposição diferenciada a variações de preços e sua interação com medidas institucionais, incorporando efeitos fixos territoriais e temporais. A interpretação causal dependerá da validade das exposições e de uma estratégia específica para a fiscalização, que pode responder à própria conversão. A heterogeneidade entre atividades e biomas será examinada conforme o suporte amostral. O estudo pretende contribuir para a compreensão da formação econômica do estoque ambiental e para a discussão de políticas de conservação sensíveis às condições territoriais e aos horizontes das decisões produtivas. Esta versão apresenta o projeto empírico para qualificação; a harmonização das bases e a construção das conversões brutas antecedem a estimação.

\medskip
\noindent\textbf{Palavras-chave:} conversão da vegetação; retorno esperado; fiscalização ambiental; uso da terra; Mato Grosso.

